
\documentclass{article}
\usepackage{colortbl}
\usepackage{makecell}
\usepackage{multirow}
\usepackage{supertabular}

\begin{document}

\newcounter{utterance}

\centering \large Interaction Transcript for game `cladder', experiment `full\_v1.5\_default', episode 525 with llama{-}3.1{-}8b{-}instruct.
\vspace{24pt}

{ \footnotesize  \setcounter{utterance}{1}
\setlength{\tabcolsep}{0pt}
\begin{supertabular}{c@{$\;$}|p{.15\linewidth}@{}p{.15\linewidth}p{.15\linewidth}p{.15\linewidth}p{.15\linewidth}p{.15\linewidth}}
    \# & \multicolumn{2}{c}{Player} && \multicolumn{2}{c}{Game Master} \\
    \hline

    \theutterance \stepcounter{utterance}  
    & & & \multicolumn{4}{p{0.6\linewidth}}{
        \cellcolor[rgb]{0.9,0.9,0.9}{
            \makecell[{{p{\linewidth}}}]{
                \texttt{\tiny{[P1$\langle$GM]}}
                \texttt{You are an expert in causal inference. The following question is not a typical commonsense query, but rather a meticulously designed question created by a professor specializing in causal inference, intended to assess the students' mastery of the course content.} \\
\texttt{Imagine a self{-}contained, hypothetical world with only the following conditions, and without any unmentioned factors or causal relationships: Unobserved confounders has a direct effect on drug taken and cholesterol level. Treatment assignment has a direct effect on drug taken. Drug taken has a direct effect on cholesterol level. Unobserved confounders is unobserved. For patients not assigned the drug treatment, the probability of low cholesterol is 63\%. For patients assigned the drug treatment, the probability of low cholesterol is 52\%. For patients not assigned the drug treatment, the probability of taking of all assigned drugs is 37\%. For patients assigned the drug treatment, the probability of taking of all assigned drugs is 74\%. Will taking of all assigned drugs increase the chance of low cholesterol?} \\
\texttt{Start your answer with 'Yes' or 'No', followed by additional reasoning or evidence to support your explanation.} \\
            }
        }
    }
     \\ \\

    \theutterance \stepcounter{utterance}  
    & \multicolumn{4}{p{0.6\linewidth}}{
        \cellcolor[rgb]{0.561,0.737,0.561}{
            \makecell[{{p{\linewidth}}}]{
                \texttt{\tiny{[P1$\rangle$GM]}}
                \texttt{No.} \\
\\ 
\texttt{To answer this question, we need to apply the principles of causal inference, specifically the concept of confounding and the causal diagram. } \\
\\ 
\texttt{Given the conditions:} \\
\\ 
\texttt{1. Unobserved confounders has a direct effect on drug taken and cholesterol level.} \\
\texttt{2. Treatment assignment has a direct effect on drug taken.} \\
\texttt{3. Drug taken has a direct effect on cholesterol level.} \\
\texttt{4. Unobserved confounders is unobserved.} \\
\\ 
\texttt{We can represent this causal diagram as follows:} \\
\\ 
\texttt{Unobserved Confounders {-}> Drug Taken} \\
\texttt{Unobserved Confounders {-}> Cholesterol Level} \\
\texttt{Treatment Assignment {-}> Drug Taken} \\
\texttt{Drug Taken {-}> Cholesterol Level} \\
\\ 
\texttt{The question asks whether taking of all assigned drugs increases the chance of low cholesterol. To answer this, we need to compare the probability of low cholesterol among patients who take all assigned drugs (T) with those who do not take all assigned drugs (NT).} \\
\\ 
\texttt{However, we cannot directly compare these probabilities because the causal diagram shows that there is a confounding variable (Unobserved Confounders) that affects both drug taken and cholesterol level. This means that the observed association between drug taken and cholesterol level may be due to the confounding variable, rather than a direct causal effect of the drug on cholesterol level.} \\
\\ 
\texttt{To address this issue, we can use the concept of counterfactuals. Let's define the following counterfactuals:} \\
\\ 
\texttt{{-} Y(T) = Cholesterol level if patient takes all assigned drugs} \\
\texttt{{-} Y(NT) = Cholesterol level if patient does not take all assigned drugs} \\
\\ 
\texttt{We are given the following probabilities:} \\
\\ 
\texttt{{-} P(Y(NT) = Low Cholesterol) = 63\% (for patients not assigned the drug treatment)} \\
\texttt{{-} P(Y(T) = Low Cholesterol) = 52\% (for patients assigned the drug treatment)} \\
\\ 
\texttt{However, we cannot directly compare these probabilities because they are based on different treatment assignments. To make a fair comparison, we need to adjust for the confounding variable.} \\
\\ 
\texttt{Using the causal diagram, we can see that the confounding variable affects both drug taken and cholesterol level. This means that the observed association between drug taken and cholesterol level is due to the confounding variable, rather than a direct causal effect of the drug on cholesterol level.} \\
\\ 
\texttt{To adjust for the confounding variable, we can use the concept of instrumental variables. However, in this case, we do not have an instrumental variable that affects treatment assignment but not the confounding variable.} \\
\\ 
\texttt{Alternatively, we can use the concept of regression adjustment. However, without additional information about the confounding variable, we cannot adjust for it.} \\
\\ 
\texttt{Given the information provided, we cannot conclude that taking of all assigned drugs increases the chance of low cholesterol. The observed association between drug taken and cholesterol level may be due to the confounding variable, rather than a direct causal effect of the drug on cholesterol level.} \\
\\ 
\texttt{Therefore, the answer is No.} \\
            }
        }
    }
    & & \\ \\

    \theutterance \stepcounter{utterance}  
    & & & \multicolumn{2}{p{0.3\linewidth}}{
        \cellcolor[rgb]{0.9,0.9,0.9}{
            \makecell[{{p{\linewidth}}}]{
                \texttt{\tiny{[GM$|$GM]}}
                \texttt{No} \\
            }
        }
    }
    & & \\ \\

    \theutterance \stepcounter{utterance}  
    & & & \multicolumn{2}{p{0.3\linewidth}}{
        \cellcolor[rgb]{0.9,0.9,0.9}{
            \makecell[{{p{\linewidth}}}]{
                \texttt{\tiny{[GM$|$GM]}}
                \texttt{no} \\
            }
        }
    }
    & & \\ \\

    \theutterance \stepcounter{utterance}  
    & & & \multicolumn{2}{p{0.3\linewidth}}{
        \cellcolor[rgb]{0.9,0.9,0.9}{
            \makecell[{{p{\linewidth}}}]{
                \texttt{\tiny{[GM$|$GM]}}
                \texttt{game\_result = WIN} \\
            }
        }
    }
    & & \\ \\

\end{supertabular}
}

\end{document}
